% !TeX program = xelatex
% ============================================================================
%  第 14 课　复习、总结与展望
% ----------------------------------------------------------------------------
%  编译        XeLaTeX，编译两遍
%  样式        ../../00-共享/preamble.tex
%  说明        本课不发单页；不学新知识，只把全课串成一条线
%  图片        一律 \imgph 占位，不要在这里手画 tikz
% ============================================================================

\documentclass[10pt]{beamer}
% ============================================================================
%  讲义样式文件 —— 字体 / 配色 / 校徽栏 / 顶部导航栏 / 页脚 / 常用方框
%  所有个性化设置都集中在这里，改一次即可用于所有讲义
%
%  用法：主文件里写
%      \documentclass[10pt]{beamer}
%      \input{preamble.tex}
%  依赖：同目录下的 logo_white.png（顶部校徽）
%  编译：必须用 XeLaTeX，而且必须编译两遍
%        （顶部导航栏的节名与小圆点记录在 .nav 文件里，只编一遍不会更新）
% ============================================================================
\usepackage{ctex}  % 支持中文
\usepackage{amsmath, amssymb}  % 数学公式支持
\usepackage{fontspec}  % 字体设置
\usepackage{unicode-math}  % 支持Unicode数学字体
\usepackage{xcolor}  % 颜色
\usepackage{tikz}  % 用于绘制页脚色条
\usepackage{booktabs}  % 表格

% 图片搜索路径 ------------------------------------------------------------------
% 校本课的资料分目录存放，本文件可能被 02-课件/第XX课/ 、03-讲义/ 下的主文件调用，
% 这里把几种可能的相对位置都列上，logo_white.png 放在 00-共享/ 即可被找到。
% 最后一项 00-共享/图片/ 是给 \imgfig 用的插图目录（图与课件不在同一层，靠它找到）。
\graphicspath{{./}{./图片/}{../../00-共享/}{../00-共享/}{00-共享/}{../../00-共享/图片/}{../00-共享/图片/}{00-共享/图片/}}

% 设置字体
\setCJKmainfont{SimSun}  % 设置中文字体为宋体
\setmainfont{Palatino Linotype}  % 设置英文字体为Palatino Linotype
\setmathfont{Latin Modern Math}  % 设置数学字体为Latin Modern Math
% 注意：beamer 默认把正文字体族设成无衬线（sans），Windows 下对应的中文字体是
% 微软雅黑，所以上面这行 SimSun 目前对正文不起作用。想让正文变成宋体，取消
% 下面这行的注释即可（标题栏、页脚、方框标题不受影响）。
% \renewcommand{\familydefault}{\rmdefault}

% 定义颜色 - 均为绿色系
\definecolor{maingreen}{RGB}{34,139,34}  % 森林绿（导航栏底色）
\definecolor{lightgreen}{RGB}{240,255,240}  % 蜜瓜白（帧标题栏底色）
\definecolor{accentgreen}{RGB}{46,139,87}  % 海绿（强调色）
\definecolor{textgreen}{RGB}{0,128,0}  % 纯绿（标题文字）
\definecolor{darkgreen}{RGB}{20,90,20}  % 深绿（顶部校徽栏底色）

% 设置Beamer主题
\usetheme{default}
\usecolortheme{default}

% 自定义主题
\setbeamercolor{background canvas}{bg=white}
\setbeamercolor{frametitle}{fg=textgreen, bg=lightgreen}
\setbeamercolor{title}{fg=textgreen}
\setbeamercolor{subtitle}{fg=textgreen}
\setbeamercolor{section in toc}{fg=accentgreen}
\setbeamercolor{normal text}{fg=black}
\setbeamercolor{structure}{fg=accentgreen}
\setbeamercolor{itemize item}{fg=accentgreen}
\setbeamercolor{itemize subitem}{fg=accentgreen}
\setbeamercolor{itemize subsubitem}{fg=accentgreen}
\setbeamercolor{block title}{fg=white, bg=accentgreen}
\setbeamercolor{block body}{bg=lightgreen}
% example 环境（底层是 exampleblock）用的是独立的一组颜色名，default 配色主题下
% 它们的底色是空的，所以只设 block 的话“例”会变成一行没有框的标题，这里接过来
\setbeamercolor{block title example}{parent=block title}
\setbeamercolor{block body example}{parent=block body}

% ==================== 常用定理类方框（五类配色）====================
% 五类各自的标题条底色，正文底色由同色自动调浅（!8!white）
\definecolor{boxThm}{RGB}{20,90,20}    % 一、定理类：深绿
\definecolor{boxDef}{RGB}{46,139,87}   % 二、定义类：海绿
\definecolor{boxExm}{RGB}{34,139,34}   % 三、例题类：森林绿
\definecolor{boxPrf}{RGB}{85,107,47}   % 四、过程类：橄榄绿
\definecolor{boxNote}{RGB}{0,121,107}  % 五、注记类：青绿

% 方框本体：先套本类配色，再调用 beamer 的 block（改动只在环境内部生效）
\newcommand{\gbox}[2]{%
  \setbeamercolor{block title}{fg=white,bg=#1}%
  \setbeamercolor{block body}{bg=#1!8!white}%
  \begin{block}{#2}}
\newcommand{\gboxend}{\end{block}}
% 标题后的可选补充名：\begin{theorem}[勾股定理] → 定理（勾股定理）
\newcommand{\gsub}[1]{\ifx\relax#1\relax\else（#1）\fi}


% 一、定理类（定理/推论/性质/命题/引理）
\renewenvironment{theorem}[1][]{\gbox{boxThm}{定理\gsub{#1}}}{\gboxend}
\renewenvironment{corollary}[1][]{\gbox{boxThm}{推论\gsub{#1}}}{\gboxend}
\newenvironment{character}[1][]{\gbox{boxThm}{性质\gsub{#1}}}{\gboxend}
\newenvironment{proposition}[1][]{\gbox{boxThm}{命题\gsub{#1}}}{\gboxend}
\renewenvironment{lemma}[1][]{\gbox{boxThm}{引理\gsub{#1}}}{\gboxend}

% 二、定义类（定义/公理）
\renewenvironment{definition}[1][]{\gbox{boxDef}{定义\gsub{#1}}}{\gboxend}
\newenvironment{axiom}[1][]{\gbox{boxDef}{公理\gsub{#1}}}{\gboxend}

% 三、例题类（例/练习/习题）
\renewenvironment{example}[1][]{\gbox{boxExm}{例\gsub{#1}}}{\gboxend}
\newenvironment{exercise}[1][]{\gbox{boxExm}{练习\gsub{#1}}}{\gboxend}
\newenvironment{homework}[1][]{\gbox{boxExm}{习题\gsub{#1}}}{\gboxend}

% 四、过程类（证明/解/法1~法4）
\renewenvironment{proof}[1][]{\gbox{boxPrf}{证明\gsub{#1}}}{\qedsymbol\gboxend}
\newenvironment{jie}[1][]{\gbox{boxPrf}{解\gsub{#1}}}{\gboxend}
\newenvironment{way1}[1][]{\gbox{boxPrf}{法1\gsub{#1}}}{\gboxend}
\newenvironment{way2}[1][]{\gbox{boxPrf}{法2\gsub{#1}}}{\gboxend}
\newenvironment{way3}[1][]{\gbox{boxPrf}{法3\gsub{#1}}}{\gboxend}
\newenvironment{way4}[1][]{\gbox{boxPrf}{法4\gsub{#1}}}{\gboxend}

% 五、注记类（经验总结/注/值得注意的是/分析/提示）
\newenvironment{experience}[1][]{\gbox{boxNote}{经验总结\gsub{#1}}}{\gboxend}
\newenvironment{zhu}[1][]{\gbox{boxNote}{注\gsub{#1}}}{\gboxend}
\newenvironment{att}[1][]{\gbox{boxNote}{值得注意的是\gsub{#1}}}{\gboxend}
\newenvironment{analyse}[1][]{\gbox{boxNote}{分析\gsub{#1}}}{\gboxend}
\newenvironment{hint}[1][]{\gbox{boxNote}{提示\gsub{#1}}}{\gboxend}

% 顶部导航栏（miniframes 主题：显示各节名称 + 帧圆点）
% 节名与圆点记录在 .nav 文件中，改动章节后需要编译两遍才会同步
\useoutertheme[subsection=false]{miniframes}
\setbeamercolor{section in head/foot}{fg=yellow, bg=maingreen}
\setbeamercolor{mini frame}{fg=yellow, bg=maingreen}
\setbeamercolor{mini frame in current subsection}{fg=yellow, bg=accentgreen}

% 顶部校徽栏：位于导航栏上方，宽度与帧标题栏一致（整页宽）
% 校训用的楷体：ctex 的 \kaishu 在本机指向未安装的 Adobe 楷体，这里直接指定系统楷体
\setCJKfamilyfont{kaiti}{KaiTi}
\setbeamercolor{logo bar}{fg=white, bg=darkgreen}
\addtobeamertemplate{headline}{%
  \begin{beamercolorbox}[wd=\paperwidth,ht=0.84cm,dp=0.20cm,leftskip=0.25cm,rightskip=0.25cm]{logo bar}%
    % -0.035cm：色条内容默认坐在基线上，会比栏的中线略高，下移这一点使其上下居中
    \raisebox{-0.035cm}{\includegraphics[height=0.72cm]{logo_white.png}}%
    % 右侧校训：字体为上面定义的楷体，也可改为 \insertsectionhead（当前节名）、\inserttitle（章节名）等
    % 0.375cm = 校徽中线（0.36cm）加汉字基线偏移，并同样下移 0.035cm 与校徽对齐
    \hfill\raisebox{\dimexpr 0.375cm - 0.5\height - 0.5\depth\relax}%
      {\CJKfamily{kaiti}\fontsize{12pt}{24pt}\selectfont 自强不息，止于至善}%
  \end{beamercolorbox}%
}{}

% % 设置边框 - 使用TikZ绘制圆角
% \setbeamertemplate{frametitle}{
%   \vspace*{0.2cm}
%   \begin{tikzpicture}
%     % 背景矩形宽度限制为标题内容宽度
%     \fill[lightgreen, rounded corners=5pt] (0,0) rectangle (\textwidth+0.6cm, 0.8cm);
%     \node[anchor=west, text width=\textwidth, xshift=0.3cm, yshift=0.4cm] {\usebeamerfont{frametitle}\insertframetitle};
%   \end{tikzpicture}
% }

% 设置项目符号
\setbeamertemplate{itemize item}{\small\raise0.5pt\hbox{\textcolor{accentgreen}{$\bullet$}}}
\setbeamertemplate{itemize subitem}{\small\raise0.3pt\hbox{\textcolor{accentgreen}{$\circ$}}}

% 设置页脚 - 使用TikZ绘制一条整页宽的色条，左侧页码、右侧姓名
\setbeamertemplate{footline}{
  \begin{tikzpicture}
    \fill[maingreen] (-1,0) rectangle (\paperwidth, 0.4cm);
    \node[yellow,anchor=east, xshift=0cm, yshift=0.2cm] {\insertframenumber/\inserttotalframenumber};
    % 原点在纸张左边缘右侧 1cm 处（色条从 x=-1cm 画起），故 \paperwidth-2cm
    % 对应纸张右边缘内缩 1cm 的位置，与左侧页码左右对称
    \node[yellow,anchor=east, xshift=\paperwidth-2cm, yshift=0.2cm] {\insertauthor};
  \end{tikzpicture}
}

% 关闭 beamer 默认的导航符号（页码与姓名已在页脚中自行绘制）
\setbeamertemplate{navigation symbols}{}

% 设置标题页面
\setbeamertemplate{title page}{
  \vspace{0.5cm}
  \begin{center}
    \begin{minipage}[c][0.68\paperheight][c]{\textwidth}
      {\usebeamerfont{title}\usebeamercolor[fg]{title}\large\inserttitle\par}
      \vspace{1.5cm}
      {\usebeamerfont{subtitle}\usebeamercolor[fg]{subtitle}\Large\insertsubtitle\par}
      \vspace{1.5cm}
      {\usebeamerfont{author}\usebeamercolor[fg]{author}\insertauthor\par}
      \vspace{0.5cm}
      {\usebeamerfont{institute}\usebeamercolor[fg]{institute}\insertinstitute\par}
      \vspace{0.5cm}
      {\usebeamerfont{date}\usebeamercolor[fg]{date}\insertdate\par}
    \end{minipage}
  \end{center}
}

% 设置块环境
\setbeamertemplate{blocks}[rounded][shadow=false]

% 设置目录
\setbeamertemplate{section in toc}{\inserttocsectionnumber.~\inserttocsection}


% ----------------------------------------------------------------------------
%  图片占位符（临时用；有图以后把那行换成 \includegraphics 即可）
%      \imgph{说明文字}          默认高 3.6cm
%      \imgph[6cm]{说明文字}     指定高度
%  会画一个绿色虚线框，框里写说明，方便对着 PDF 找位置。
\newcommand{\imgph}[2][3.6cm]{%
  \begin{center}%
  \begin{tikzpicture}%
    \node[draw=accentgreen, densely dashed, line width=1pt, rounded corners=4pt,
          fill=lightgreen, inner sep=8pt, align=center, text width=0.72\textwidth,
          minimum height=#1, font=\small]
      {\textcolor{accentgreen}{\bfseries 【图片占位】}\\[4pt] #2};%
  \end{tikzpicture}%
  \end{center}%
}

% ----------------------------------------------------------------------------
%  插图（图已就位后用这个；图放在 00-共享\图片\ 下，由 \graphicspath 找到）
%      \imgfig{文件名}                 默认宽 0.72\textwidth
%      \imgfig[0.55\textwidth]{文件名}  指定宽度
%      \imgfig[8cm]{文件名}            也可以直接给尺寸
%  图源与编译方法见 00-共享\图片\README.md。每张图是一页 PDF，只含图、无页边。
\newcommand{\imgfig}[2][0.72\textwidth]{%
  \begin{center}%
    \includegraphics[width=#1]{#2}%
  \end{center}%
}

% ----------------------------------------------------------------------------
%  可用方框环境速查（五类配色，均不编号，均支持 \begin{xxx}[补充名]）
%    定理类（深绿）  theorem 定理 / corollary 推论 / character 性质 /
%                    proposition 命题 / lemma 引理
%    定义类（海绿）  definition 定义 / axiom 公理
%    例题类（森林绿）example 例 / exercise 练习 / homework 习题
%    过程类（橄榄绿）proof 证明 / jie 解 / way1 法1 / way2 法2 /
%                    way3 法3 / way4 法4
%    注记类（青绿）  experience 经验总结 / zhu 注 / att 值得注意的是 /
%                    analyse 分析 / hint 提示
%  另有 beamer 自带的普通方框：\begin{block}{标题} ... \end{block}
% ----------------------------------------------------------------------------


% ==================== 每份课件只改这一段 ====================
\title{从数数开始的抽象代数生活}
\subtitle{第 14 课　复习、总结与展望}
\author{姜博瀚}
\institute{厦门大学附属实验中学}
\date{\today}
% ============================================================

\begin{document}

\begin{frame}
  \titlepage
\end{frame}

% ============================================================================
\section{把 14 节课串成一条线}

\begin{frame}{今天不学新东西}
  \begin{center}
    {\large 把 14 节课串成一条线。}
  \end{center}
  \pause
  \begin{center}
    {\large 所有内容你都见过，只是散着放。}\\[0.9em]
    \begin{tabular}{@{}r@{\quad}l@{}}
      \textcolor{textgreen}{\bfseries 左边} & 数着数着，就绕回来了 \\
      \textcolor{textgreen}{\bfseries 右边} & 位置换一换，还是那几个样子 \\
      \textcolor{textgreen}{\bfseries 下方} & 幂的尾巴，也绕圈 \\
    \end{tabular}
  \end{center}
\end{frame}

\begin{frame}{左边：数着数着，就绕回来了}
  \begin{itemize}
    \item 模 \(12\) 加法：\(8+7=3\)，走 \(7\) 格，绕回来了；
    \item 单位元 \(0\)：加上它等于没加；
    \item 逆元：\(1\) 与 \(11\) 配对、\(2\) 与 \(10\) 配对、\(6\) 与自己配对；
    \item 由 \(3\) 生成的 \(\{0,3,6,9\}\)：只走 \(4\) 步就回到 \(0\)。
  \end{itemize}
  \pause
  \begin{center}
    {\large 数没有变，变的只是「绕圈」的走法。}
  \end{center}
\end{frame}

\begin{frame}{右边：位置换一换，还是那几个样子}
  \begin{itemize}
    \item 正三角形：\(6\) 个对称动作，记作 \(S_3\)；
    \item 正方形：\(8\) 个对称动作，记作 \(D_4\)；
    \item 把复合当作运算：先 \(r\) 后 \(f\)，与先 \(f\) 后 \(r\) 不一样，\(rf\neq fr\)；
    \item 阶与子群：小世界嵌在大世界里，子群的大小整除群的大小。
  \end{itemize}
  \pause
  \begin{center}
    {\large 图形没有变，只是「换个位置摆放」。}
  \end{center}
\end{frame}

\begin{frame}{下方：幂的尾巴，也绕圈}
  \begin{itemize}
    \item \(2^n\) 除以 \(7\)，余数在 \(2,4,1\) 里打转；
    \item 试过许多小质数（也叫素数），归纳出猜想：\(a^{\,p-1}\equiv1\pmod p\)；
    \item 用「余数是重排」把它证出来——费马小定理；
    \item 回头回答第 1 课：\(10^{6}\equiv1\pmod 7\)，所以 \textbf{\(7\) 整除 \(999999\)}。
  \end{itemize}
\end{frame}

\begin{frame}{三条线汇到一句}
  \begin{definition}[群]
    一个集合配上一种运算，如果满足下面四条，就叫\textbf{群}：
    \begin{enumerate}
      \item \textbf{封闭}：任取 \(a,b\)，\(a*b\) 还在这里面；
      \item \textbf{结合}：\((a*b)*c=a*(b*c)\)；
      \item \textbf{单位元}：存在 \(e\)，使 \(a*e=e*a=a\)；
      \item \textbf{逆元}：每个 \(a\) 都有 \(b\)，使 \(a*b=b*a=e\)。
    \end{enumerate}
  \end{definition}
  \pause
  \begin{center}
    {\large 数、动作、置换——只要满足这四条，数学家就把它们归到一起研究。}
  \end{center}
  \pause
  \begin{att}
    这里的 \(*\) 只是个占位的记号：换成 \(+\)、\(\times\)，或者「两个动作接着做」，都行。
  \end{att}
\end{frame}

% ============================================================================
\section{抽查}

\begin{frame}{用手指回答}
  \begin{itemize}
    \item 模 \(12\) 加法里，\(5\) 的逆元是谁？
    \item 正三角形有 \(6\) 个对称动作，那正方形有几个？
    \item 三角形的对称里，\(rf\) 和 \(fr\) 一样吗？
    \item \(2^{100}\) 除以 \(7\) 余几？
  \end{itemize}
\end{frame}

\begin{frame}{刚才那四问的答案}
  \begin{itemize}
    \item \(5\) 的逆元是 \(7\)：因为 \(5+7=12\equiv0\)；
    \item 正方形有 \(8\) 个对称动作（正三角形是 \(6\) 个）；
    \item \(rf\) 和 \(fr\) 不一样，\(rf\neq fr\)；
    \item \(2^{100}\) 除以 \(7\) 余 \(2\)：\(100=6\times16+4\)，所以
      \((2^{6})^{16}\times2^{4}\equiv1\times16\equiv2\)。
  \end{itemize}
\end{frame}

\begin{frame}{对着主线图说一句话}
  \begin{itemize}
    \item 从黑板上任选一个点，说出它是什么意思，再举一个例子；
    \item 说不出来就先听别人说，不用慌——今天不考试。
  \end{itemize}
  \pause
  \begin{center}
    {\large 整门课，你只需要记住一件事：这些例子满足同样的四条规则。}
  \end{center}
\end{frame}

% ============================================================================
\section{前面还有什么}

\begin{frame}{加法与乘法，放一起}
  \footnotesize
  \setlength{\parskip}{0pt}
  \setlength{\itemsep}{1pt}
  \noindent 分配律你们早就在用了：\(3\times(4+5)=3\times4+3\times5\)。
  \begin{definition}[环]
    一个集合上同时有加法和乘法，如果下面三条都满足，就叫做\textbf{环}：
    \begin{enumerate}
      \item \textbf{加法}：满足第 4 课那四条规则（封闭、结合、单位元、逆元），而且加法可以交换：\(a+b=b+a\)；
      \item \textbf{乘法}：封闭、结合、\textbf{有单位元 \(1\)}（\(a\cdot1=1\cdot a=a\)）、而且乘法可以交换：\(a\cdot b=b\cdot a\)——但乘法不要求每个数都有逆元；
      \item \textbf{乘法对加法有分配律}：\(a\cdot(b+c)=a\cdot b+a\cdot c\)，\((b+c)\cdot a=b\cdot a+c\cdot a\)。
    \end{enumerate}
    乘法可以交换，所以它又叫\textbf{交换环}——乘法不交换的环也有，这门课不去碰它。
  \end{definition}
  \pause
  \begin{definition}[域]
    在上面这个环的基础上，如果每个非零元素都有乘法逆元（能做除法），就叫做\textbf{域}。
  \end{definition}
  \pause
  \begin{center}
    \small 整数 \(\mathbb{Z}\) 是环但不是域（\(2\) 没有整数逆元，\(\tfrac12\notin\mathbb{Z}\)）；有理数 \(\mathbb{Q}\) 是域。
  \end{center}
\end{frame}

\begin{frame}{二次方程：配方就能解}
  \small
  \setlength{\abovedisplayskip}{1pt plus 1pt}
  \setlength{\belowdisplayskip}{1pt plus 1pt}
  \setlength{\abovedisplayshortskip}{1pt}
  \setlength{\belowdisplayshortskip}{1pt}
  \begin{example}[一元二次方程]
    解 \(ax^2+bx+c=0\quad(a\neq0)\)：两边除以 \(a\)，再令中间变量
    \(t=x+\frac{b}{2a}\)，就化成
    \[ t^2=\frac{b^2-4ac}{4a^2}. \]
    令 \(\Delta=b^2-4ac\)（叫\textbf{判别式}），于是
    \[ t=\pm\frac{\sqrt{\Delta}}{2a},\qquad x=\frac{-b\pm\sqrt{b^2-4ac}}{2a}. \]
  \end{example}
  \begin{center}
    \small\textbf{一句话}：先把 \(x\) 平移到 \(t\)，方程就只剩一个平方。\\
    \small \(\Delta>0\) 两个根，\(\Delta=0\) 一个根，\(\Delta<0\) 没有实数根。\\
    \scriptsize（\(\sqrt{\ }\) 是平方根，\(\sqrt[3]{\ }\) 是立方根；这里只要认得这个记号就行。）\\
    \scriptsize（只看个样子，不用记。真正把这几个公式一条条全写出来，能占满一整块黑板；所以数学家后来才问：五次能不能也有这样一条路？）
  \end{center}
\end{frame}

\begin{frame}{三次方程：卡尔达诺设 \(t=u+v\)}
  \small
  \setlength{\abovedisplayskip}{1pt plus 1pt}
  \setlength{\belowdisplayskip}{1pt plus 1pt}
  \setlength{\abovedisplayshortskip}{1pt}
  \setlength{\belowdisplayshortskip}{1pt}
  \begin{example}[一元三次方程]
    \(t^3+pt+q=0\)（任何三次方程都能先平移消掉平方项，变成这个「缺项」的样子）。
    令
    \[ t=u+v,\qquad\text{并约定}\ uv=-\frac{p}{3}, \]
    代入后方程变成 \(u^3+v^3=-q\)，而由约定又有
    \[ u^3v^3=-\frac{p^3}{27}. \]
    所以 \(u^3\) 与 \(v^3\) 是一元二次方程
    \[ z^2+qz-\frac{p^3}{27}=0 \]
    的两个根。\textbf{解出这个二次方程，就得到 \(u^3,v^3\)，再开立方就得到 \(t=u+v\)。}
  \end{example}
  \begin{center}
    \small\textbf{一句话}：三次方程被「两个中间变量」降成了一个二次方程。\\
    \scriptsize（只看个样子，不用记。真正把这几个公式一条条全写出来，能占满一整块黑板；所以数学家后来才问：五次能不能也有这样一条路？）
  \end{center}
\end{frame}

\begin{frame}{四次方程：费拉里配两次方}
  \small
  \setlength{\abovedisplayskip}{1pt plus 1pt}
  \setlength{\belowdisplayskip}{1pt plus 1pt}
  \setlength{\abovedisplayshortskip}{1pt}
  \setlength{\belowdisplayshortskip}{1pt}
  \begin{example}[一元四次方程]
    \(t^4+pt^2+qt+r=0\)（同样先平移消掉三次项）。第一步配方：
    \[ \left(t^2+\frac{p}{2}\right)^2=-\frac{p^2}{4}+r-qt. \]
    第二步，引入中间变量 \(y\)，给左边添上 \(2y\left(t^2+\frac p2\right)+y^2\)，让它仍然是完全平方：
    \[ \left(t^2+\frac{p}{2}+y\right)^2
       =2y\,t^2-qt+\left(y^2+py-\frac{p^2}{4}+r\right). \]
    只要右边的判别式为 \(0\)（即 \(q^2=8y^3+8py^2-8ry\)，这是一个关于 \(y\) 的三次方程），
    右边也是完全平方；于是两边开方，就得到两个二次方程。
  \end{example}
  \begin{center}
    \small\textbf{一句话}：四次方程靠「中间变量 \(y\)」拆成两个二次方程——而 \(y\) 又要靠上面那个三次方程求出来。\\
    \scriptsize（只看个样子，不用记。真正把这几个公式一条条全写出来，能占满一整块黑板；所以数学家后来才问：五次能不能也有这样一条路？）
  \end{center}
\end{frame}

\begin{frame}{五次方程的故事}
  \begin{itemize}
    \item 一次、二次方程：有求根公式；
    \item 三次、四次方程：十六世纪也找到了公式；
    \item 五次方程：数学家找了三百多年。
  \end{itemize}
  \pause
  \begin{center}
    {\large 十九世纪，伽罗瓦证明：五次方程\textbf{没有}求根公式。}
  \end{center}
  \pause
  \begin{center}
    {\large 他用的工具，正是这门课讲的——\textbf{群}。\quad 他去世时只有二十岁。}
  \end{center}
\end{frame}

\begin{frame}{黑板上的最后一行}
  \begin{center}
    \vspace{0.5cm}
    {\Large 群不是新东西，}\\[0.5em]
    {\Large 它就是「绕圈」这件事的名字。}\\[1.2em]
    \pause
    {\Large\textcolor{textgreen}{\bfseries 而它，解开了三百年的谜。}}
  \end{center}
\end{frame}

% ---------------- 结尾页 ----------------
\begin{frame}
  \begin{center}
    \vspace{1.6cm}
    {\Large\textcolor{textgreen}{\bfseries 全课终}}

    \vspace{1.0cm}
    {\large\textcolor{accentgreen}{谢谢这一学期，一起去绕圈。}}
  \end{center}
\end{frame}

\end{document}